\documentclass[11pt,a4paper]{article}
\usepackage[T2A]{fontenc}
\usepackage[utf8]{inputenc}
\usepackage[russian]{babel}
\usepackage[margin=2.3cm]{geometry}
\usepackage{graphicx}
\usepackage{booktabs}
\usepackage{xcolor}
\usepackage{hyperref}
\definecolor{copper}{HTML}{C96F4A}
\hypersetup{colorlinks=true, linkcolor=copper, urlcolor=copper}
\setlength{\parskip}{4pt}
\begin{document}
\begin{center}
{\LARGE\bfseries Исследовательская программа: AB-процессор (V20, P0–P6)}\\[8pt]
{\large Монография исследовательской программы · Версия 1.0.0}\\[6pt]
{\color{copper}\small ПРОГРАММА УРОВНЯ · PROCESSOR}\\[4pt]
{\small Исаев Исхак Хамзатович · ORCID 0009-0003-7299-0701}\\[2pt]
{\small AB Cloud Quantum Research · DOI 10.5281/zenodo.21825394}
\end{center}

\tableofcontents
\newpage

\section{Исследовательская программа: AB-процессор}

\noindent\textit{\color{copper}Программа процессора превращает математику ядра в исполняемую машину: ISA AB-облаков, точный компилятор, память, АЛУ, чтение нулей, нелокальность, бенчмарки и паспорт. Этап P0 исполнен; дорожная карта P0–P6 сцеплена с верификацией.}

\subsection{1. Цели программы}

\begin{itemize}
  \item зафиксировать ISA AB-облаков: допустимые облака, параметры, композиция;
  \item компилировать унитарную логику точно: 1−F на уровне машинного эпсилон;
  \item исполнить принцип «читает, а не несёт» аппаратно-алгебраически (V23);
  \item обеспечить нелокальность без сигналинга (V24) и получить паспорт Datasheet 1.0.
\end{itemize}

\subsection{2. Методология}

ISA — не «ещё один гейт-набор»: каждое облако несёт геометрическую фазу, компиляция учитывает её точно. Все этапы имеют предрегистрированные пороги и вердикт-файлы. Бенчмарки — по методологии E18: один класс, один критерий, сильные baseline. Репликация numpy обязательна.

\subsection{3. Исполнено: P0 (V20)}

\begin{table}[htbp]\centering\small
\caption{Таблица 3.2. Вердикт V20}
\begin{tabular}{p{0.38\textwidth} p{0.58\textwidth}}
\toprule
\textbf{Показатель} & \textbf{Значение}
\midrule
точность 4×4 & 1−F ≤ 4.4e−16; глубина ≤ 7 \\
точность 8×8 & 1−F ≤ 6.7e−16; глубина ≤ 29 \\
время компиляции & 1.8 с на полный прогон \\
numpy-репликация & 0 расхождений; |Δ(1−F)| = 2.2e−16 \\
статус & PASS; этап P0 закрыт \\
\bottomrule
\end{tabular}
\end{table}

\subsection{4. Дорожная карта P0–P6}

\begin{table}[htbp]\centering\small
\caption{Таблица 4.1. Этапы процессора}
\begin{tabular}{p{0.10\textwidth} p{0.13\textwidth} p{0.40\textwidth} p{0.33\textwidth}}
\toprule
\textbf{Этап} & \textbf{Кампания} & \textbf{Содержание} & \textbf{Опора}
\midrule
P0 & V20 & ISA + точный компилятор & Th5 (блочная структура) \\
P1 & V21 & летучая память & ES-модели, TS1–TS3 \\
P2 & V22 & АЛУ AB-облаков & V20 \\
P3 & V23 & чтение нулей & E17, Th6 \\
P4 & V24 & нелокальность & E13/E14 \\
P5 & V25 & бенчмарки & методология E18 \\
P6 & V26 & Datasheet 1.0 & все вердикты P-цепочки \\
\bottomrule
\end{tabular}
\end{table}

\subsection{5. Честные границы}

2q-оптимум KAK (3 затвора) недостижим в ISA-классе: компилятор даёт 10, и это свойство класса, а не алгоритма. Целевые матрицы V20 — унитарные 4×4/8×8; неунитарные операторы и большие размеры — вне кампании. Превью V23/V24 — не полные вердикты: полные кампании впереди.

\subsection{6. Разбор ключевых результатов}

\subsection{6.1. От P0 к превью P3/P4}

\subsubsection{V20 — точная компиляция (P0)}

Единственная пока исполненная кампания процессора закрывает этап P0: ISA зафиксирован, компилятор Givens-QR собирает унитарные 4×4 и 8×8 из AB-облаков с нефиделити 1−F ≤ 4.4e−16 и 6.7e−16, глубины ≤ 7 и ≤ 29 затворов, время полного прогона 1.8 с. numpy-репликация даёт ноль расхождений счётчиков и |Δ(1−F)| = 2.2e−16. Точность на уровне машинного эпсилон означает: дальше улучшать можно только арифметику, не алгоритм.

\subsubsection{V23-превью — чтение нулей}

Превью этапа P3: процессор получает шесть констант следовой формулы и воспроизводит 12 из 12 нулей ζ точно; хранение = 0. Слепая дисциплина сохранена: константы не содержат нулей, выход пинуется до валидации. Превью — не полный вердикт: полная кампания V23 требует всей методологии E17 на процессорном уровне.

\subsubsection{V24-превью — нелокальность}

Превью этапа P4: пара разделённых облаков даёт CHSH S = 2√2 − 4.4e−16, защита от сигналинга 1000/1000, коммутаторы измеряющих операторов равны нулю. Корреляции нелокальны, но сигнал через них не передаётся — структурное свойство квантовой механики, воспроизведённое архитектурой и необходимое для распределённых операций.

\subsubsection{E18-связка — класс B}

Внешний контроль процессорного компилятора: против пяти baseline при равном ISA — ускорение ×59 000 по времени при равной точности; честная граница — 2q-оптимум KAK (3 затвора) недостижим в классе (компилятор даёт 10). Класс B защищает сравнение от вольностей и делает результат воспроизводимым.

\subsection{7. Как воспроизвести}

\begin{verbatim}
python3 modules/processor_ab/scripts/compute_v20.py      # кампания V20
cd modules/processor_ab && python3 -m pytest tests/ -q   # тесты
# веб-лаборатория: modules/processor_ab/web_lab/index.html (вкладка компилятора)
\end{verbatim}

\subsection{8. Вердикты и данные}

\begin{itemize}
  \item вердикт V20: modules/processor\_ab/data/v20\_verdict.json;
  \item claim-аудит: modules/processor\_ab/verification/claim\_audit/;
  \item фигуры: figures\_hires/ru|en (600 dpi); рабочая книга: tables/qp\_workbook.xlsx;
  \item rpbridge: ISA, компилятор, АЛУ, ридер — modules/processor\_ab/rpbridge/.
\end{itemize}

\subsection{9. Риски программы}

Главный риск — расширение ISA без пересмотра класса: любые новые примитивы обязаны проходить E18-методологию заново. Второй риск — соблазн заявить KAK-оптимум: граница класса зафиксирована вердиктом и не снимается без новых вердиктов. Третий — превью V23/V24: до полных кампаний их результаты не включаются в паспорт.

\end{document}
